\documentclass[11pt]{article}
\usepackage[margin=1.1in]{geometry}
\usepackage{amsmath,amssymb,amsthm}
\usepackage{enumitem}
\usepackage{booktabs}
\usepackage[colorlinks=true,linkcolor=blue,citecolor=blue,urlcolor=blue]{hyperref}

\newtheorem{theorem}{Theorem}[section]
\newtheorem{lemma}[theorem]{Lemma}
\newtheorem{proposition}[theorem]{Proposition}
\newtheorem{corollary}[theorem]{Corollary}
\newtheorem{remark}[theorem]{Remark}
\theoremstyle{definition}
\newtheorem{definition}[theorem]{Definition}
\newtheorem{objection}{Objection}[section]

\newcommand{\pP}{\textbf{[P]}}
\newcommand{\pC}{\textbf{[C]}}
\newcommand{\pO}{\textbf{[O]}}
\newcommand{\Tr}{\operatorname{Tr}}

\title{\textbf{Emergent Gravity and a New Lagrangian:\\
A Theory of Everything We Know}}
\author{Robert W.\ McGwier\\
\small CTO and Staff Scientist, Cohere Technology Group, Herndon, VA}
\date{August 7, 2026 --- Version 1}

\begin{document}
\maketitle

\begin{abstract}
\noindent
We consolidate a research program \cite{McGwier2026a,McGwier2026b,
McGwier2026c} into a single document with a single thesis: the
gravitational sector of known physics is entanglement bookkeeping, and the
unique effective Lagrangian this bookkeeping selects, once the empirically
confirmed matter content is installed, is the Standard Model minimally
coupled to Einstein--Cartan gravity with cosmological constant. Part~I
establishes the mechanism: in holographic conformal field theories the
first law of entanglement, imposed on all ball-shaped regions, is
equivalent to the linearized Einstein equations; the superposition
principle for geometry follows as a theorem; and the extension to
Riemann--Cartan geometry is derived, with torsion produced by the same
bookkeeping wherever the state carries spin density. Part~II assembles the
Lagrangian and its certifications, including the single parameter-free
prediction beyond known physics: a universal gravitational-strength
axial--axial contact interaction among all fermions. The title is exact
and limited: this is a theory of everything \emph{we know}---curvature,
inertia, torsion, and geometric superposition derived; gauge structure,
masses, and generations installed as empirical input; finality expressly
disclaimed. Appendices~A and~B establish the domain of validity (AdS
versus de Sitter) and the theorem-grade status of Wick rotation.
Appendix~C answers, in numbered form, the objections a skeptical referee
should raise. Every material claim is labeled \pP{} (proved/derived),
\pC{} (conjectured with support), or \pO{} (open).
\end{abstract}

\section*{Epistemic Conventions}
Every material claim is tagged: \pP{} proved or derived within a stated
framework, with the derivation published in the primary literature cited
at the point of use or presented in full here; \pC{} conjectured,
supported by nontrivial consistency checks but lacking proof; \pO{} open.
No claim is asserted without a tag or citation. Numerical agreement is not
proof. Consensus is not proof.

\part*{Part I: Gravity from Entanglement}

\section{The Mechanism}

\subsection{Kinematics}

\begin{definition}
For a spatial region $A$ and global state $|\Psi\rangle$, the reduced
state is $\rho_A=\Tr_{\bar A}|\Psi\rangle\langle\Psi|$, the modular
Hamiltonian $H_A=-\log\rho_A$ (up to normalization), and the entanglement
entropy $S_A=-\Tr\rho_A\log\rho_A$.
\end{definition}

Two exact results anchor everything. For the Rindler wedge in any Wightman
QFT, the vacuum modular Hamiltonian is $2\pi$ times the boost generator
(Bisognano--Wichmann \cite{BisognanoWichmann1976}). For a ball $B$ of
radius $R$ in a CFT vacuum (Casini--Huerta--Myers \cite{CHM2011}),
\begin{equation}
H_B=2\pi\!\int_B d^{d-1}x\;\frac{R^2-|\vec x-\vec x_0|^2}{2R}\,
T_{00}(x)\,.
\label{eq:CHM}
\end{equation}

\begin{theorem}[First law of entanglement]
\label{thm:firstlaw}
For any family $\rho_A(\lambda)$ with $\rho_A(0)=\rho_A^{\mathrm{vac}}$,
\begin{equation}
\delta S_A=\delta\langle H_A\rangle=\Tr\!\big(\delta\rho_A H_A\big)
\quad\text{at first order.}
\label{eq:firstlaw}
\end{equation}
\end{theorem}

\begin{proof}
$\delta S_A=-\Tr(\delta\rho_A\log\rho_A^{\mathrm{vac}})-\Tr\,\delta\rho_A$;
the second term vanishes by trace preservation, the first is
$\delta\langle H_A\rangle$. \pP{}
\end{proof}

The right-hand side of \eqref{eq:firstlaw}---the \emph{modular energy}
$E_{\mathrm{mod}}[\delta\rho]=\Tr(\delta\rho_A H_A)$---is the linear
functional on which the entire program runs (Appendix~C, Objection~C.3).
The exact identity behind it is
$S(\rho_A\|\rho_A^{\mathrm{vac}})=\Delta\langle H_A\rangle-\Delta S_A
\ge0$ \cite{Casini2008,BlancoCasini2013}: entropy's nonlinearity is
quarantined in the relative entropy, which is second order and positive.

\subsection{Dynamics: the equivalence theorem}

\begin{theorem}[Faulkner--Guica--Hartman--Myers--Van Raamsdonk
\cite{Faulkner2014}; see also \cite{Lashkari2014}]
\label{thm:fghmv}
Let a holographic CFT$_d$ state
$|\Psi\rangle=|0\rangle+\lambda|\delta\Psi\rangle$ be dual to
$g=g_{\mathrm{AdS}}+\lambda h$, with $\delta S_B$ computed gravitationally
by Ryu--Takayanagi \cite{RyuTakayanagi2006,HRT2007} and
$\delta\langle H_B\rangle$ by \eqref{eq:CHM}. Then
$\delta S_B=\delta\langle H_B\rangle$ for all balls of all radii, centers,
and frames \emph{if and only if} $h_{\mu\nu}$ satisfies the Einstein
equations linearized about AdS, sourced by
$\delta\langle T_{\mu\nu}\rangle$. \pP{}
\end{theorem}

\begin{remark}[The mechanism, stated exactly]
\label{rem:mechanism}
An excitation---a quasiparticle---curves spacetime through the
perturbation it induces in the \emph{vacuum} entanglement across every
partition, quantified by \eqref{eq:firstlaw}; consistency of these
$\delta S_B$ across all partitions forces an emergent metric obeying
linearized Einstein dynamics. Intrinsic entanglement between excitations
is a microstate correlation, not a source: two excitations in a Bell state
versus a product state with identical stress profiles produce the
identical linearized metric. \pP{} (Appendix~C, Objections C.1--C.2.)
\end{remark}

Beyond linear order: relative-entropy identities reproduce the nonlinear
Einstein equations through second order \cite{Faulkner2017,FLM2013};
Jacobson's entanglement equilibrium \cite{Jacobson2016} proposes the full
nonlinear equations for general backgrounds \pC{}; connectivity from
entanglement \cite{VanRaamsdonk2010,MaldacenaSusskind2013} and tensor
networks \cite{Swingle2012} supply the structural picture \pP{}/\pC{} as
itemized in \cite{McGwier2026a}.

\section{Superposition as a Theorem}
\label{sec:superposition}

\begin{proposition}[Linear order]
Both sides of \eqref{eq:firstlaw} are linear in $\delta\rho$, so the map
$\delta\rho\mapsto h$ of Theorem~\ref{thm:fghmv} is linear:
superposed source perturbations map to superposed metric perturbations,
coherences included. Weak-field superposition of gravity is a corollary,
not a postulate; naive semiclassical sourcing by
$\langle T_{\mu\nu}\rangle$ of a macroscopic superposition is excluded
within the framework and experimentally \cite{PageGeilker1981}. \pP{}
\end{proposition}

\begin{proposition}[Macroscopic superpositions]
In holographic quantum error correction the Ryu--Takayanagi area is an
operator in the center of the region algebra
\cite{ADH2015,Harlow2017,HaPPY2015}:
$S(\rho_A)=\Tr(\rho\,\widehat{\mathrm{Area}}/4G_N)+S_{\mathrm{alg}}$.
Superpositions of macroscopically distinct geometries are superpositions
of area eigenstates with Born-rule statistics; the center supplies the
approximate superselection by which classical geometry decoheres. \pP{}
within AdS/CFT; beyond \pO{}. No collapse mechanism is required
\cite{Penrose1996}; gravitationally mediated entanglement experiments
\cite{Bose2017,MarlettoVedral2017} discriminate empirically. \pC{}
pending data.
\end{proposition}

\section{GR plus Torsion, Derived}
\label{sec:torsion}

In first-order variables ($e^a$, $\omega^{ab}$; torsion
$T^a=de^a+\omega^a{}_b\wedge e^b$), the Palatini action
\begin{equation}
S=\frac{1}{4\kappa}\int\epsilon_{abcd}\,e^a\wedge e^b\wedge R^{cd}
+S_m[e,\omega,\psi]\,,\qquad\kappa=8\pi G\,,
\label{eq:ECaction}
\end{equation}
yields under independent variation
\cite{Kibble1961,Sciama1964,HehlRMP1976,Trautman2006}
\begin{equation}
\tfrac12\epsilon_{abcd}e^b\wedge R^{cd}=\kappa\tau_a\,,\qquad
T^\lambda{}_{\mu\nu}+\delta^\lambda_\mu T^\sigma{}_{\nu\sigma}
-\delta^\lambda_\nu T^\sigma{}_{\mu\sigma}=\kappa s^\lambda{}_{\mu\nu}\,.
\label{eq:ECeqs}
\end{equation}
The Cartan equation is algebraic: torsion is locked pointwise to spin
density, does not propagate, and vanishes in vacuum. Eliminating it yields
Riemannian GR with Belinfante source plus the Hehl--Datta contact term
\cite{HehlDatta1971}; for Dirac matter,
$\mathcal L_{\mathrm{HD}}=-\tfrac{3\kappa}{16}
(\bar\psi\gamma^5\gamma^\mu\psi)(\bar\psi\gamma^5\gamma_\mu\psi)$. \pP{}
The thermodynamic derivation extends to this setting
\cite{DLP2017,DDS2018}, with entropy $=\mathrm{Area}/4G$ robust in
first-order variables \cite{JacobsonMohd2015}. \pP{} within assumptions.

\begin{proposition}[Torsionful first law, linear order]
\label{prop:ECfirstlaw}
Linearizing about a maximally symmetric vacuum: background torsion
vanishes (Lorentz invariance + algebraicity); the system decomposes into a
metric sector obeying linearized Einstein equations with Belinfante source
and an algebraic torsion sector fixed by
$\delta\langle s^\lambda{}_{\mu\nu}\rangle$; Hehl--Datta is quadratic and
drops. Since the CHM Hamiltonian \eqref{eq:CHM} is built from the
Belinfante tensor, Theorem~\ref{thm:fghmv} extends verbatim: the first
law is equivalent to linearized Einstein--Cartan gravity, with torsion
produced by the same bookkeeping wherever the state carries spin. \pP{}
\end{proposition}

The companion note \cite{McGwier2026b} resolves the case of explicit
spin-current sources: for minimal Dirac matter only the axial (totally
antisymmetric) source couples, pure-gauge sources act only through their
curl, the torsionful first law holds unconditionally on null planes (via
the Markov property \cite{CasiniTesteTorroba2017} and the FLPW deformation
theory \cite{FLPW2016,SarosiUgajin2018}) and conditionally for finite
balls; mixed-symmetry sources through non-minimal matter are relevant
deformations that destroy the construction. \pP{}/\pO{} as itemized there.
This \emph{modular selection rule} is what Part~II uses.

\part*{Part II: The Lagrangian}

\section{The New Standard Model}
\label{sec:nsm}

\begin{definition}
The New Standard Model (NSM) is the action \eqref{eq:NSM}---the full
$SU(3)_c\times SU(2)_L\times U(1)_Y$ Standard Model
\cite{Glashow1961,Weinberg1967,Salam1968}, three generations, Higgs sector
\cite{ATLAS2012,CMS2012}, neutrino masses \cite{SuperK1998} (Dirac or
Weinberg-operator, empirically open \pO{}), minimally coupled through
vielbein and independent spin connection to Einstein--Cartan gravity with
cosmological constant---\emph{together with} the certifications of
Section~\ref{sec:certifications}.
\end{definition}

\begin{equation}
\boxed{\;
S_{\mathrm{NSM}}
=\frac{1}{4\kappa}\!\int\!\epsilon_{abcd}\Big(e^a\wedge e^b\wedge R^{cd}
-\frac{\Lambda}{6}\,e^a\wedge e^b\wedge e^c\wedge e^d\Big)
+\int d^4x\;e\,\mathcal{L}_{\mathrm{SM}}[e,\omega,\psi,A,H]\;}
\label{eq:NSM}
\end{equation}
with $\mathcal L_{\mathrm{SM}}$ the standard gauge, fermion, Higgs, and
Yukawa terms, fermions covariantized with the full connection
$D_\mu=\partial_\mu+\tfrac14\omega_{ab\mu}\gamma^{ab}+i g\!\cdot\!A_\mu$.
Gauge field strengths use the exterior derivative alone (gauge invariance
forbids torsion coupling \cite{HehlRMP1976,Shapiro2002}); scalars carry no
spin. \emph{Only fermions talk to torsion.} \pP{}

Eliminating the algebraic torsion yields the single interaction beyond the
SM and Einstein gravity:
\begin{equation}
\boxed{\;
\mathcal L_{\mathrm{HD}}=-\frac{3\kappa}{16}\,
\mathcal J_5^\mu\mathcal J_{5\mu}\,,\qquad
\mathcal J_5^\mu=\!\!\sum_{f\in\mathrm{all\ fermions}}\!\!
\bar\psi_f\gamma^5\gamma^\mu\psi_f\;}
\label{eq:HD}
\end{equation}
---universal, flavor-blind, cross-species, gravitational strength, with
\emph{no free parameter}: the coefficient is fixed by $G$. An output, not
an input. \pP{}

\section{The Certifications}
\label{sec:certifications}

\begin{theorem}[Gravity certified]
At linear order the metric dynamics of \eqref{eq:NSM} is equivalent to the
entanglement first law on all balls (Theorem~\ref{thm:fghmv},
Proposition~\ref{prop:ECfirstlaw}): curvature and inertia are the unique
consistent response to perturbations of vacuum entanglement. \pP{} within
the holographic regime; nonlinear torsionful equations under Clausius
assumptions \cite{Jacobson1995,DLP2017}. 
\end{theorem}

\begin{theorem}[Torsion structure certified]
The minimal axial torsion coupling of \eqref{eq:NSM} is the unique
structure surviving the modular selection rule of \cite{McGwier2026b}:
mixed-symmetry couplings through non-minimal matter are relevant
deformations destroying the modular framework. \pP{} at the level of the
dimensional and symmetry argument.
\end{theorem}

\begin{theorem}[Superposition certified]
Section~\ref{sec:superposition} applies verbatim: the NSM requires no
metric quantization as an independent postulate and no collapse
mechanism. \pP{} within AdS/CFT; \pO{} beyond.
\end{theorem}

\begin{remark}
No analogous certification exists for the gauge and Yukawa structure:
the matter content is empirical input, and whether gauge dynamics,
chirality, and three generations admit an entanglement-theoretic
derivation is the deepest question this program can pose. \pO{}
\end{remark}

\section{Falsifiable Consequences; Declared Limits}

Consequences: (1) the parameter-free contact term \eqref{eq:HD},
consistent with torsion bounds \cite{KRT2008,Shapiro2002}, undetectable at
collider energies \pP{}; (2) torsion-induced bounce replacing the initial
singularity \cite{Trautman1973,Poplawski2010} \pC{}; (3) gravitationally
mediated entanglement as a live falsifier of the gravitational sector
\cite{Bose2017,MarlettoVedral2017} \pC{}; (4) infrared deviations at the
Milgrom scale if the de Sitter extension follows the emergent route
\cite{Verlinde2017,Milgrom1983} \pC{}.

Declared limits: the NSM does not unify gauge couplings, explain masses or
generations, resolve strong CP, supply a dark-matter particle, determine
$\Lambda$, or provide its own UV completion---the last by construction,
since the Lagrangian is asserted to be emergent bookkeeping for a deeper
description \pO{}. \emph{A theory of everything we know is not a theory of
everything}; on this program's epistemology, finality is not claimable.

\appendix

\section{Domain of Validity: AdS versus de Sitter}
\label{app:adsds}

\subsection{Why the proofs live in AdS}
The \pP{}-grade chain requires three ingredients AdS supplies and de
Sitter does not: a timelike boundary carrying a complete unitary CFT with
the Maldacena dictionary \cite{Maldacena1998} (de Sitter's candidate
dS/CFT \cite{Strominger2001} is Euclidean, generically non-unitary, and
\pO{}); a global timelike Killing structure guaranteeing the
Bisognano--Wichmann machinery (de Sitter has only observer-dependent
static patches with Gibbons--Hawking horizon entropy
\cite{GibbonsHawking1977}, microscopic description \pO{}); and
theorem-grade entanglement technology
\cite{RyuTakayanagi2006,FLM2013,EngelhardtWall2015}, unavailable at that
grade in de Sitter \pO{}. Rule: certify in AdS, extrapolate with labels.

\subsection{Why the mechanism governs our space}
\begin{proposition}[Flat-space inputs]
The quantum-informational inputs---CHM \eqref{eq:CHM}, the first law,
relative-entropy positivity and monotonicity, the null-plane Markov
property \cite{CasiniTesteTorroba2017}---are theorems of QFT on Minkowski
space, independent of holography. The vacuum of the Standard-Model fields
in our universe carries the structure the mechanism runs on. \pP{}
\end{proposition}
\begin{proposition}[Locality]
The thermodynamic route \cite{Jacobson1995,DLP2017} uses only local
Rindler horizons, available through every point of every spacetime, and
yields the field equations pointwise; AdS/CFT upgrades the local inputs to
theorems. \pP{} structure; general-background transfer
\cite{Jacobson2016} \pC{}.
\end{proposition}
\begin{proposition}[Infrared decoupling of $\Lambda$]
With $\Lambda\simeq1.1\times10^{-52}\,\mathrm{m}^{-2}$, the local
relevance measure $\Lambda L^2$ is $\sim10^{-52}$ (laboratory),
$\sim10^{-30}$ (1 AU), $\sim10^{-11}$ (10 kpc), $\mathcal O(1)$ only at
the Hubble radius. Below the Hubble scale the AdS/dS/flat distinction is
invisible to local experiment; the certified regime is
$\Lambda L^2\ll1$. \pP{}
\end{proposition}
\begin{proposition}[Where the difference bites]
The sign of $\Lambda$ is structural only near the horizon scale, where de
Sitter has an observer horizon with entropy \cite{GibbonsHawking1977} and
possibly a volume-law entanglement contribution \cite{Verlinde2017}. Any
modification appears at accelerations $\sim cH_0$; Milgrom's
$a_0\simeq1.2\times10^{-10}\,\mathrm{m\,s^{-2}}\approx cH_0/6$
\cite{Milgrom1983} sits exactly there. Coincidence \pP{}; explanation
\pC{}; microtheory \pO{}. The proof's edge coincides with observation's
anomalies.
\end{proposition}

\section{The Status of Wick Rotation}
\label{app:wick}

Euclidean continuation enters at three load-bearing points, each covered
by a theorem. \textbf{Reconstruction:} Osterwalder--Schrader
\cite{OS1973} uniquely reconstructs a unitary, causal, positive-energy
Lorentzian QFT from reflection-positive Euclidean data; the boundary
theories used here satisfy the hypotheses. \pP{}
\textbf{Instantiation:} Bisognano--Wichmann
\cite{BisognanoWichmann1976} \emph{is} the rotation, proved once at the
foundational level: wedge modular flow is the boost, with
\eqref{eq:CHM} its conformal image; Tomita--Takesaki theory
\cite{Haag1996} performs no rotation at all; the Unruh temperature
\cite{Unruh1976} is the KMS countersignature. \pP{}
\textbf{The fragile uses are absent from \pP{} claims:} (i) the replica
limit $n\to1$ \cite{LM2013} is never performed at linear order---the first
law is an operator identity and Theorem~\ref{thm:fghmv} uses RT as an
independently established input; (ii) the Euclidean gravitational path
integral, unbounded below \cite{GHP1978}, is never performed---gravity is
emergent and every gravitational statement is boundary entanglement;
(iii) de Sitter Euclidean vacua are genuinely subtle and lie entirely
inside the \pO{} region of Appendix~\ref{app:adsds}---the two frontiers
coincide, a consistency check on the labeling. \pP{} Fermion
$\gamma^5$ conventions touch one certified claim (the anomaly caveat in
\cite{McGwier2026b}) and are fixed convention-independently by the index
theorem \cite{Fujikawa1979}. \pP{}

\section{Anticipated Objections, Answered}
\label{app:objections}

Each objection below was raised, in some form, during the development of
this program. Answers cite the theorem that settles the point or the tag
that concedes it.

\begin{objection}[``This is Verlinde's entropic gravity, and/or it runs
on `relative entropy.'\,'']
Neither. The entropy in Theorems~\ref{thm:firstlaw}--\ref{thm:fghmv} is
\emph{entanglement entropy across a partition}---relationally
\emph{defined} (no partition, no entropy) but not ``relational entropy''
as a term of art, and not Verlinde's coarse-grained screen entropy
\cite{Verlinde2011}, which requires $N\gg1$ bits at finite temperature and
is undefined for a single excitation. Relative entropy enters in one exact
role: as the positive second-order remainder
$S(\rho\|\rho_{\mathrm{vac}})=\Delta\langle H\rangle-\Delta S$
\cite{Casini2008,BlancoCasini2013}, whose positivity delivers energy
conditions and second-order dynamics \cite{Faulkner2017}. Verlinde's 2016
volume-law proposal \cite{Verlinde2017} is a \pC{} extension confined to
Appendix~\ref{app:adsds}'s frontier, never load-bearing for a \pP{}. \pP{}
\end{objection}

\begin{objection}[``So entangled particles gravitate toward each other
because they're entangled?'']
No, and the theorem's content forbids it. Curvature at linear order is
fixed entirely by $\delta\langle T_{\mu\nu}\rangle$
(Remark~\ref{rem:mechanism}): a Bell pair and a product pair with
identical stress profiles produce identical metrics. The gravitating
object is the excitation's perturbation of \emph{vacuum} entanglement,
ball by ball; pair entanglement is a microstate correlation invisible to
single-region reduced states. Whether ER=EPR
\cite{MaldacenaSusskind2013} endows pairs with a Planckian geometric dual
is \pC{} and carries no classical curvature signature in any case. \pP{}
\end{objection}

\begin{objection}[``A thermodynamic origin for gravity forfeits quantum
superposition. What is even linear here?'']
The linear functional is the \emph{modular energy}
$E_{\mathrm{mod}}[\delta\rho]=\Tr(\delta\rho_A H_A)$, and the first law
says $\delta S_A$ \emph{equals} it: entropy's nonlinearity is entirely
quarantined in the positive second-order relative entropy. Both sides of
the equivalence in Theorem~\ref{thm:fghmv} are therefore linear in
$\delta\rho$, making weak-field superposition a corollary
(Section~\ref{sec:superposition}); macroscopic superpositions of distinct
geometries are area-operator superpositions \cite{Harlow2017}. The naive
semiclassical alternative is excluded experimentally
\cite{PageGeilker1981}; the collapse alternative \cite{Penrose1996} is
empirically discriminable \cite{Bose2017,MarlettoVedral2017}. \pP{}
within stated frameworks.
\end{objection}

\begin{objection}[``Torsion is ad hoc, dead since the 1970s, and
unobservable.'']
Torsion is not added; it is what the first-order formalism \emph{already
contains} the moment fermions exist: spinors couple to the connection, the
connection's variation is the Cartan equation \eqref{eq:ECeqs}, and
refusing torsion is the ad hoc step. It is algebraic (no new propagating
fields), experimentally bounded and consistent \cite{KRT2008,Shapiro2002},
selected uniquely by the modular argument \cite{McGwier2026b}, and its one
consequence \eqref{eq:HD} is parameter-free. ``Unobservable at colliders''
is conceded \pP{} and is exactly what a gravitational-strength contact
term must be; the high-density regime \cite{Trautman1973,Poplawski2010}
is where it earns its keep. \pC{} there.
\end{objection}

\begin{objection}[``Why is everything proved in anti-de Sitter when the
universe is de Sitter?'']
Because in AdS the claim is a theorem and in de Sitter it is a hope, and
the discipline of this document is never to blur that line. The full
argument is Appendix~\ref{app:adsds}: AdS supplies the boundary, the
dictionary, and the theorem-grade technology; the mechanism so certified
is \emph{local}, built from flat-space theorems our fields satisfy; the
measured $\Lambda$ decouples below the Hubble scale
($\Lambda L^2\ll1$); and the one regime where the certification runs
out is the one regime where observation shows anomalies. AdS is the wind
tunnel, not the sky. \pP{}/\pC{}/\pO{} as tagged there.
\end{objection}

\begin{objection}[``Euclidean methods smuggle in unjustified analytic
continuation.'']
Appendix~\ref{app:wick}: the rotation used is licensed by
Osterwalder--Schrader, instantiated by Bisognano--Wichmann,
countersigned by the Unruh/KMS thermality, and the three genuinely fragile
continuations---replica limits, the Euclidean gravitational path integral,
de Sitter vacua---are respectively never used at linear order,
structurally bypassed by emergence, and coincident with the declared
\pO{} frontier. \pP{}
\end{objection}

\begin{objection}[``You have a Lagrangian for QFT plus torsionful
gravity---so you claim a fundamental theory.'']
Two levels, kept separate. At the effective level, yes:
\eqref{eq:NSM}--\eqref{eq:HD} is the complete EFT of known physics,
quantizable below $M_{\mathrm{Pl}}$, and \emph{certified} in its
gravitational and torsion structure. At the fundamental level, no, by
construction: in this framework the Lagrangian is emergent bookkeeping;
the microscopic description is the boundary theory of a holographic pair,
and no first-principles selection of that pair exists \pO{}. The sharpest
open prize is whether the Hehl--Datta coefficient itself can be derived
from relative-entropy positivity at joint second order, which would make
the entire EFT an entanglement theorem. \pO{}
\end{objection}

\begin{objection}[``The title claims a theory of everything.'']
Read it in full: a theory of everything \emph{we know}. Derived:
curvature, inertia, torsion, geometric superposition, and one new
interaction with no free parameter. Installed as empirical input: gauge
group, representations, Yukawas, generations, $\Lambda$. Declared open:
their origin, the UV completion, de Sitter. The program's own
epistemology---emergence without a final layer---forbids the stronger
reading, and Section~7 enumerates the exclusions with the same rigor as
the claims. A referee who finds the title too strong is invited to find a
single sentence in the body that is. \pP{} (as a statement about the
document).
\end{objection}

\section*{Certified Summary}
Gravity is entanglement bookkeeping \pP{}; its torsionful extension is
derived and modular-selected \pP{}; geometry superposes as a theorem
\pP{}; the unique resulting effective Lagrangian of known physics is
\eqref{eq:NSM} with the single parameter-free consequence \eqref{eq:HD}
\pP{}; the domain of validity and the analytic-continuation machinery are
theorem-anchored with their frontiers labeled \pP{}/\pO{}; and everything
not derived is named. This is the current certified layer---the next
turtle, load-bearing and inspected.

\section*{Acknowledgments}
The author acknowledges the use of Claude (Anthropic, claude.ai) for
assistance in drafting, literature verification, and \LaTeX{} preparation.
All technical claims and epistemic classifications were reviewed and
verified by the author, who bears sole responsibility for the content.

\begin{thebibliography}{99}\small

\bibitem{Glashow1961}
S.~L.~Glashow, Nucl.\ Phys.\ \textbf{22}, 579 (1961).
\bibitem{Kibble1961}
T.~W.~B.~Kibble, J.\ Math.\ Phys.\ \textbf{2}, 212 (1961).
\bibitem{Sciama1964}
D.~W.~Sciama, Rev.\ Mod.\ Phys.\ \textbf{36}, 463 (1964).
\bibitem{Weinberg1967}
S.~Weinberg, Phys.\ Rev.\ Lett.\ \textbf{19}, 1264 (1967).
\bibitem{Salam1968}
A.~Salam, in \emph{Elementary Particle Theory}, ed.\ N.~Svartholm
(Almqvist \& Wiksell, 1968), p.~367.
\bibitem{HehlDatta1971}
F.~W.~Hehl and B.~K.~Datta, J.\ Math.\ Phys.\ \textbf{12}, 1334 (1971).
\bibitem{OS1973}
K.~Osterwalder and R.~Schrader, Commun.\ Math.\ Phys.\ \textbf{31}, 83
(1973); II, \textbf{42}, 281 (1975).
\bibitem{Trautman1973}
A.~Trautman, Nature Phys.\ Sci.\ \textbf{242}, 7 (1973).
\bibitem{Hawking1975}
S.~W.~Hawking, Commun.\ Math.\ Phys.\ \textbf{43}, 199 (1975).
\bibitem{BisognanoWichmann1976}
J.~J.~Bisognano and E.~H.~Wichmann, J.\ Math.\ Phys.\ \textbf{17}, 303
(1976).
\bibitem{HehlRMP1976}
F.~W.~Hehl, P.~von der Heyde, G.~D.~Kerlick and J.~M.~Nester,
Rev.\ Mod.\ Phys.\ \textbf{48}, 393 (1976).
\bibitem{Unruh1976}
W.~G.~Unruh, Phys.\ Rev.\ D \textbf{14}, 870 (1976).
\bibitem{GibbonsHawking1977}
G.~W.~Gibbons and S.~W.~Hawking, Phys.\ Rev.\ D \textbf{15}, 2738 (1977).
\bibitem{GHP1978}
G.~W.~Gibbons, S.~W.~Hawking and M.~J.~Perry, Nucl.\ Phys.\ B
\textbf{138}, 141 (1978).
\bibitem{Fujikawa1979}
K.~Fujikawa, Phys.\ Rev.\ Lett.\ \textbf{42}, 1195 (1979).
\bibitem{PageGeilker1981}
D.~N.~Page and C.~D.~Geilker, Phys.\ Rev.\ Lett.\ \textbf{47}, 979
(1981).
\bibitem{Milgrom1983}
M.~Milgrom, Astrophys.\ J.\ \textbf{270}, 365 (1983).
\bibitem{Jacobson1995}
T.~Jacobson, Phys.\ Rev.\ Lett.\ \textbf{75}, 1260 (1995).
\bibitem{Haag1996}
R.~Haag, \emph{Local Quantum Physics}, 2nd ed.\ (Springer, 1996).
\bibitem{Penrose1996}
R.~Penrose, Gen.\ Rel.\ Grav.\ \textbf{28}, 581 (1996).
\bibitem{Maldacena1998}
J.~M.~Maldacena, Adv.\ Theor.\ Math.\ Phys.\ \textbf{2}, 231 (1998).
\bibitem{SuperK1998}
Super-Kamiokande Collaboration, Phys.\ Rev.\ Lett.\ \textbf{81}, 1562
(1998).
\bibitem{Strominger2001}
A.~Strominger, JHEP \textbf{10}, 034 (2001).
\bibitem{Shapiro2002}
I.~L.~Shapiro, Phys.\ Rept.\ \textbf{357}, 113 (2002).
\bibitem{Trautman2006}
A.~Trautman, in \emph{Encyclopedia of Mathematical Physics} (Elsevier,
2006), arXiv:gr-qc/0606062.
\bibitem{RyuTakayanagi2006}
S.~Ryu and T.~Takayanagi, Phys.\ Rev.\ Lett.\ \textbf{96}, 181602 (2006).
\bibitem{HRT2007}
V.~E.~Hubeny, M.~Rangamani and T.~Takayanagi, JHEP \textbf{07}, 062
(2007).
\bibitem{Casini2008}
H.~Casini, Class.\ Quantum Grav.\ \textbf{25}, 205021 (2008).
\bibitem{KRT2008}
V.~A.~Kosteleck\'y, N.~Russell and J.~D.~Tasson, Phys.\ Rev.\ Lett.\
\textbf{100}, 111102 (2008).
\bibitem{VanRaamsdonk2010}
M.~Van Raamsdonk, Gen.\ Rel.\ Grav.\ \textbf{42}, 2323 (2010).
\bibitem{Poplawski2010}
N.~J.~Pop\l awski, Phys.\ Lett.\ B \textbf{694}, 181 (2010).
\bibitem{Verlinde2011}
E.~Verlinde, JHEP \textbf{04}, 029 (2011).
\bibitem{CHM2011}
H.~Casini, M.~Huerta and R.~C.~Myers, JHEP \textbf{05}, 036 (2011).
\bibitem{ATLAS2012}
ATLAS Collaboration, Phys.\ Lett.\ B \textbf{716}, 1 (2012).
\bibitem{CMS2012}
CMS Collaboration, Phys.\ Lett.\ B \textbf{716}, 30 (2012).
\bibitem{Swingle2012}
B.~Swingle, Phys.\ Rev.\ D \textbf{86}, 065007 (2012).
\bibitem{MaldacenaSusskind2013}
J.~Maldacena and L.~Susskind, Fortsch.\ Phys.\ \textbf{61}, 781 (2013).
\bibitem{BlancoCasini2013}
D.~D.~Blanco, H.~Casini, L.-Y.~Hung and R.~C.~Myers, JHEP \textbf{08},
060 (2013).
\bibitem{LM2013}
A.~Lewkowycz and J.~Maldacena, JHEP \textbf{08}, 090 (2013).
\bibitem{FLM2013}
T.~Faulkner, A.~Lewkowycz and J.~Maldacena, JHEP \textbf{11}, 074 (2013).
\bibitem{Lashkari2014}
N.~Lashkari, M.~B.~McDermott and M.~Van Raamsdonk, JHEP \textbf{04}, 195
(2014).
\bibitem{Faulkner2014}
T.~Faulkner, M.~Guica, T.~Hartman, R.~C.~Myers and M.~Van Raamsdonk,
JHEP \textbf{03}, 051 (2014).
\bibitem{ADH2015}
A.~Almheiri, X.~Dong and D.~Harlow, JHEP \textbf{04}, 163 (2015).
\bibitem{HaPPY2015}
F.~Pastawski, B.~Yoshida, D.~Harlow and J.~Preskill, JHEP \textbf{06},
149 (2015).
\bibitem{EngelhardtWall2015}
N.~Engelhardt and A.~C.~Wall, JHEP \textbf{01}, 073 (2015).
\bibitem{JacobsonMohd2015}
T.~Jacobson and A.~Mohd, Phys.\ Rev.\ D \textbf{92}, 124010 (2015).
\bibitem{Jacobson2016}
T.~Jacobson, Phys.\ Rev.\ Lett.\ \textbf{116}, 201101 (2016).
\bibitem{FLPW2016}
T.~Faulkner, R.~G.~Leigh, O.~Parrikar and H.~Wang, JHEP \textbf{09}, 038
(2016).
\bibitem{Harlow2017}
D.~Harlow, Commun.\ Math.\ Phys.\ \textbf{354}, 865 (2017).
\bibitem{Bose2017}
S.~Bose et al., Phys.\ Rev.\ Lett.\ \textbf{119}, 240401 (2017).
\bibitem{MarlettoVedral2017}
C.~Marletto and V.~Vedral, Phys.\ Rev.\ Lett.\ \textbf{119}, 240402
(2017).
\bibitem{Verlinde2017}
E.~Verlinde, SciPost Phys.\ \textbf{2}, 016 (2017).
\bibitem{DLP2017}
R.~Dey, S.~Liberati and D.~Pranzetti, Phys.\ Rev.\ D \textbf{96}, 124032
(2017).
\bibitem{CasiniTesteTorroba2017}
H.~Casini, E.~Test\'e and G.~Torroba, J.\ Phys.\ A \textbf{50}, 364001
(2017).
\bibitem{Faulkner2017}
T.~Faulkner, F.~M.~Haehl, E.~Hijano, O.~Parrikar, C.~Rabideau and
M.~Van Raamsdonk, JHEP \textbf{08}, 057 (2017).
\bibitem{SarosiUgajin2018}
G.~S\'arosi and T.~Ugajin, JHEP \textbf{01}, 012 (2018).
\bibitem{DDS2018}
T.~De Lorenzo, E.~De Paoli and S.~Speziale, Phys.\ Rev.\ D \textbf{98},
064053 (2018).
\bibitem{McGwier2026a}
R.~W.~McGwier, ``Entanglement as the origin of spacetime curvature,''
v4, Zenodo (2026), DOI:
\href{https://doi.org/10.5281/zenodo.21821890}{10.5281/zenodo.21821890}.
\bibitem{McGwier2026b}
R.~W.~McGwier, ``Spin-current sources and ball modular Hamiltonians,''
Zenodo (2026), DOI:
\href{https://doi.org/10.5281/zenodo.21824934}{10.5281/zenodo.21824934}.
\bibitem{McGwier2026c}
R.~W.~McGwier, ``The New Standard Model for unified physics,''
Zenodo (2026), DOI:
\href{https://doi.org/10.5281/zenodo.21829536}{10.5281/zenodo.21829536}.

\bibitem{McGwier2026d}
R.~W.~McGwier, ``Condition NL for finite balls: a spectral reduction, the unconditional metric sector, and a universal kinematic residue,''
Zenodo (2026), DOI:
\href{https://doi.org/10.5281/zenodo.21836572}{10.5281/zenodo.21836572}.

\bibitem{McGwier2026e}
R.~W.~McGwier, ``The Hehl--Datta coefficient from relative-entropy positivity,''
Zenodo (2026), DOI:
\href{https://doi.org/10.5281/zenodo.21836757}{10.5281/zenodo.21836757}.

\end{thebibliography}

\end{document}
